\documentclass[11pt]{amsart}
\usepackage[T1]{fontenc}
\usepackage{lmodern}
\usepackage{amsmath,amssymb,amsthm}
\usepackage[numbers,sort]{natbib}
\usepackage[hidelinks]{hyperref}
\usepackage{bm}

\newtheorem{theorem}{Theorem}
\newtheorem{lemma}[theorem]{Lemma}
\newtheorem{corollary}[theorem]{Corollary}
\newtheorem{proposition}[theorem]{Proposition}
\newtheorem{conjecture}[theorem]{Conjecture}
\theoremstyle{definition}
\newtheorem{definition}[theorem]{Definition}

\newcommand{\Ter}{\mathcal S}
\newcommand{\red}{\twoheadrightarrow}
\newcommand{\hred}{\twoheadrightarrow_h}
\newcommand{\conv}{=_\beta}
\newcommand{\sz}[1]{\lvert#1\rvert}
\newcommand{\rest}{\mathrm{rest}}
\newcommand{\nf}{\mathrm{nf}}
\newcommand{\fold}{\mathrm{fold}}
\newcommand{\pop}{\mathrm{pop}}
\newcommand{\Hc}{\mathcal H}
\newcommand{\Ic}{\mathcal I}
\newcommand{\Gc}{\mathcal G}
\newcommand{\Kc}{\mathcal K}
\newcommand{\Cc}{\mathcal C}
\newcommand{\Self}{\mathrm{SELF}}
\newcommand{\sqs}{\sqsubseteq}
\newcommand{\psq}{\sqsubset}

\title[Counter Runs Are Determined by Their Teeth]{Counter Runs of the Combinator $\bm{S}$ Are Determined by Their Teeth}
\author{Claude Opus \& Fable with N.\ E.\ Davis}
\address{ZeroAttest, Champaign, Illinois, USA}
\email{NDavis@ZeroAttest.com}
\urladdr{https://orcid.org/0000-0002-1131-8319}
\date{\today}

\begin{document}
\begin{abstract}
The head run of a closed term built from the combinator $S$, if it never reaches a head normal
form, splits off an infinite sequence of arguments, its teeth. Convertible terms of this kind have
convertible teeth, and the converse, CT, is open. We prove CT against every term $p$ whose run is
a counter run: from some point on its states lie in an explicit grammar of counters, and its teeth
normalize. Every term without a head normal form whose teeth are convertible with those of $p$ is
then convertible with $p$. This covers the runs of $U_j\,U_m$ and of $X\,(S\,X\,c)$. The proof
follows the two runs in normal form. In the grammar, moves that copy the whole heap into the stack
recur, and opposite such a move a second run with the same teeth reaches the same state, at once
or after a few pops. The proof is not mechanized, except for two finite facts about a 38-element
model of head normal forms, which are checked by kernel evaluation in Lean~4.
\end{abstract}

\keywords{combinatory logic, combinator $S$, head normal form, normal form, convertibility}
\maketitle

\section{Introduction}

$\Ter$ is the set of closed terms built from $S$ by application, with $S\,x\,y\,z\to x\,z\,(y\,z)$ at
any position; $\red$ is the reflexive--transitive closure of $\to$, and $\conv$ is convertibility.
The \emph{size} $\sz t$ is the number of occurrences of $S$ in $t$. A step raises the size by
$\sz z-1\ge0$, so $t\red u$ implies $\sz t\le\sz u$; we cite this as (Size). A term
\emph{normalizes} if it reduces to a normal form, which is then unique, $\nf(t)$, by
Church--Rosser (CR)~\citep{terese}.

We use the terminology of~\citep[\S\S1--2]{s11}. Every $t\in\Ter$ is $S\,u_1\cdots u_m$ for a
unique $m\ge0$; the $u_i$ are its \emph{items}, and $t$ is a \emph{head normal form} (hnf) iff
$m\le2$. For $m\ge3$ the \emph{head step} is $t\to u_1\,u_3\,(u_2\,u_3)\,u_4\cdots u_m$; head steps
give the \emph{head run} of $t$, whose terms are its \emph{frames}, and $\hred$ is their closure.
The term $t$ \emph{has an hnf} if its head run reaches one, and $t$ is \emph{literal} if it is its
own hnf. A frame $f\,w$ is \emph{stable} if $f$ has no hnf. If $t$ has no hnf, its head run
reaches a first stable frame $f\,w$, and $\tau_1(t)=w$, $\rest(t)=f$; the \emph{teeth} of $t$ are
$\tau_{n+1}(t)=\tau_1(\rest^n(t))$, and \emph{move} $n$ goes from $\rest^{n-1}(t)$ to $\rest^n(t)$
and produces $\tau_n(t)$. Convertible terms without an hnf have convertible teeth
\citep[Lemma~5]{s11}. The converse is open:

\begin{conjecture}[CT]
If $p,q\in\Ter$ have no hnf and $\tau_n(p)\conv\tau_n(q)$ for all $n\ge1$, then $p\conv q$.
\end{conjecture}

\noindent The companion paper~\citep{s11} proves CT when $p$ is one explicit term. Here we prove it
for every $p$ whose run reaches the following grammar and whose teeth normalize. Write
$T=S\,S\,S$, $X=S\,T\,S$, $U=S\,T\,T$, $D(v)=S\,v\,(S\,v)$, and let $c_0=S$, $c_{j+1}=S\,c_j$ be
the \emph{combs}, $\Cc=\{c_j:j\ge0\}$.

\begin{definition}\label{def:grammar}
$\Hc$ and $\Ic$ are the least sets of normal terms such that
\begin{itemize}
\item[(H0)] $X,U\in\Hc$; \quad (H1) $H\in\Hc\Rightarrow S\,H\in\Hc$;
\item[(H2)] $P\in\Hc$ has two items and $c$ is normal $\Rightarrow S\,P\,c\in\Hc$;
\item[(H3)] $P\in\Hc$ has one item and $Q\in\Hc\cup\Ic\cup\{S\}\Rightarrow S\,P\,Q\in\Hc$;
\item[(H4)] $P'\in\Hc$ has one item and $c\in\Cc\Rightarrow S\,(S\,P')\,c\in\Hc$;
\item[(I)] $Q\in\Hc$ and $c\in\Cc\Rightarrow S\,c\,Q\in\Ic$.
\end{itemize}
$\Hc_2$ is the set of members of $\Hc$ with two items, and $\Gc=\Hc_2\times\Hc$.
\end{definition}

\begin{theorem}\label{thm:main}
Let $(Z,W)\in\Gc$ be such that every tooth of $Z\,W$ normalizes. Let $p,q\in\Ter$ have no hnf,
with $\rest^e(p)\conv Z\,W$ for some $e\ge0$ and $\tau_n(q)\conv\tau_n(p)$ for all $n\ge1$. Then
$q\conv p$.
\end{theorem}

\noindent So CT holds against $Z\,W$, against its states, which are convertible with members of
$\Gc$ (Corollary~\ref{cor:recur}), and against every $p$ whose run in normal form
(Section~\ref{sec:pairs}) reaches $\Gc$ and has normalizing teeth; we call such runs
\emph{counter runs}. Two families satisfy the hypothesis. Let
$U_0=U$, $U_{m+1}=D(U_m)$, and let $\Kc$ be the least set of normal terms that contains every
normal term with at most one item and contains $S\,S\,b$ whenever it contains $b$.

\begin{corollary}\label{cor:U}
For $j,m\ge0$ every tooth of $U_j\,U_m$ normalizes, and CT holds against every $p$ without an hnf
such that $\rest^e(p)\conv U_j\,U_m$ for some $e\ge0$.
\end{corollary}

\begin{corollary}\label{cor:X}
Let $Y=X$, or $Y=S\,X\,c$ with $c$ normal. Every tooth of $X\,Y$ normalizes iff $Y=X$ or $c\in\Kc$;
then CT holds against every $p$ without an hnf such that $\rest^e(p)\conv X\,Y$ for some $e\ge0$.
\end{corollary}

\emph{Idea.} We follow both runs in normal form: from some point on each state is
convertible with $A\,W$, a \emph{normal pair} of a two-item normal stack $A$ and a normal heap $W$
(Lemmas~\ref{lem:pair} and~\ref{lem:reach}). In $\Gc$ the run never leaves $\Gc$ (Theorem~\ref{thm:closure}), and
infinitely often the stack $X$ or $U$ copies a two-item heap $W$ into the stack, with tooth
$S\,W\,(S\,W)$ or $S\,W\,(S\,W\,(S\,W))$ (Corollary~\ref{cor:recur}). Opposite such a move the
other run's move is pinned down by inverting normal forms (Lemmas~\ref{lem:fa} and~\ref{lem:tee}),
and it either reaches the same state or stands where a later tooth must differ
(Lemmas~\ref{lem:dh} and~\ref{lem:alpha}). Hnf-lessness of members applied to members comes from
two finite facts about the block model of~\citep{s1} (Lemma~\ref{lem:tables}).

\emph{Assumptions.} The proof uses (CR), standardization (Std) and the existence of stable frames
(F2) through Lemmas~3--7 of~\citep{s11}, which states them and cites their proofs; the cited proof of
(F2) uses the block model of~\citep{s1}. The block model enters otherwise only through
Lemma~\ref{lem:tables}: the folds of $X$ and $U$, the characterization (M2) of hnf-less terms by
their block, and two table facts, T16 and T17, checked by kernel evaluation in Lean. These facts and
(CR), (Std), (F2) are mechanized; the proofs in this paper are not. The proof case-splits on whether a term has an hnf, in~\citep[Lemma~4]{s11} and in
the normal-form computation of a move (Lemma~\ref{lem:pair}); (M2) decides that property. The
other steps are inductions on sizes and depths and case distinctions on the number of items of
normal terms; we have not audited them for other classical principles.

\section{Normal pairs}\label{sec:pairs}

We cite Lemmas~3 (head runs), 4 (moves), 5 (convertible terms), 6 (finite meeting) and 7 (right
suffixes) of~\citep{s11} as L3--L7. By L4, a move of $f=f^-\,u$ without an hnf is of kind (A),
where $f^-$ has no hnf, $\tau_1(f)=u$ and $\rest(f)=f^-$; or of kind (B): $f^-\hred S\,a_0\,b_0$,
$w_0=u$, $w_{i+1}=b_i\,w_i$, $a_i\,w_i\hred S\,a_{i+1}\,b_{i+1}$ for $i<k$, $a_k\,w_k$ has no hnf,
$\tau_1(f)=b_k\,w_k$ and $\rest(f)=a_k\,w_k$. \emph{Entry} $i$ is \emph{literal} if
$S\,a_i\,b_i$ is reached in no head step, that is, $f^-=S\,a_0\,b_0$ ($i=0$) or
$a_{i-1}\,w_{i-1}=S\,a_i\,b_i$ ($i\ge1$). If it is not, the last frame before $S\,a_i\,b_i$ is
$S\,S\,y\,a_i$ and $b_i=y\,a_i$ for a term $y$ (L3(a)).

A right context of depth $d$ is $x_1\,(x_2\,(\cdots(x_d\,[\,])\cdots))$; $e$ is a right suffix of $t$
at depth $d$ if $t$ is such a context filled with $e$. We write $e\sqs t$, and $e\psq t$ if
$d\ge1$. Right suffixes of $t$ at different depths have different sizes, and $\sqs$ is
transitive. Since a subterm of a normal term is normal, L7(a), applied to $t\red\nf(t)$, gives:
\begin{itemize}
\item[(RS)] if $e\psq t$ and $t$ normalizes, then $e$ normalizes and $\nf(e)\psq\nf(t)$.
\end{itemize}

\begin{lemma}\label{lem:sub}
If $t$ normalizes, every subterm of $t$ normalizes, and $t$ has an hnf.
\end{lemma}

\begin{proof}
Induction on the length of a reduction $t\red n$ to normal form, for all subterm occurrences at
once. If $t=n$, its subterms are normal. Otherwise $t=C[S\,x\,y\,z]\to t_1=C[x\,z\,(y\,z)]$, and
the subterms of $t_1$ normalize. A subterm $u$ of $t$ disjoint from the redex or inside $x$, $y$ or
$z$ occurs in $t_1$; if $u$ contains the redex, $u\to u_1$ with $u_1$ a subterm of $t_1$. Otherwise
$u$ is $S$, $S\,x$ or $S\,x\,y$, which normalize since $x$ and $y$ do. A normal term has at most two
items, so $t$ reduces to an hnf and has one by L3(c).
\end{proof}

So a term with a subterm without an hnf does not normalize, nor does any term that reduces to one
(CR). A \emph{normal pair} $(A,W)$ consists of normal terms such that $A\,W$ has no hnf; then $A$
has two items, $A=S\,A^0\,B^0$, since otherwise $A\,W$ would be an hnf. The teeth and moves of
$(A,W)$ are those of $A\,W$; $A$ is its \emph{stack} and $W$ its \emph{heap}.

\begin{lemma}\label{lem:pair}
Let $(A,W)$ be a normal pair whose first tooth normalizes. Its move is of kind \textup{(B)}, with
$a_0=A^0$, $b_0=B^0$, $w_0=W$, entry $0$ literal, and all $a_i$, $b_i$, $w_i$ \textup{(}$i\le k$\textup{)}
normalize. With $A^i=\nf(a_i)$, $B^i=\nf(b_i)$, $W^i=\nf(w_i)$: $A^i\,W^i$ has an hnf iff $i<k$, and
then $\nf(A^i\,W^i)=S\,A^{i+1}\,B^{i+1}$ and $W^{i+1}=\nf(B^i\,W^i)$. The pair $(A^k,W^k)$ is a normal
pair, $\rest(A\,W)\conv A^k\,W^k$, and the tooth has the normal form $N=\nf(B^k\,W^k)$. Moreover:
\begin{itemize}
\item[(a)] $W^i\psq W^{i+1}$ for $i<k$, and $W^k\psq N$;
\item[(b)] if entry $i\ge1$ is literal, $B^i=W^{i-1}$; if it is not, $B^i\ne S$;
\item[(c)] if $k=0$, the new state is $(A^0,W)$, and $A^0$ is a proper subterm of $A$.
\end{itemize}
\end{lemma}

\begin{proof}
$A$ is literal with two items, so the move is of kind (B) with these data. The tooth $b_k\,w_k$
normalizes, and $w_k=b_{k-1}\,(\cdots(b_0\,w_0))$ contains $b_i$, $w_i$ for $i<k$; so all of them
normalize (Lemma~\ref{lem:sub}). If some entry is not literal, let $i$ be the last one; then
$b_i=y\,a_i$ and $a_i=S^{k-i}\,a_k$, so $a_k$ is a subterm of $b_i$, which is $b_k$ or a subterm of
$w_k$. Otherwise $a_k$ is a subterm of $A^0$. Either way $a_k$ normalizes; and for $i<k$,
$a_i\,w_i\red S\,a_{i+1}\,b_{i+1}$, so by downward induction $a_i\,w_i$ and $a_i$ normalize. As
$a_i\,w_i\conv A^i\,W^i$, the two have an hnf together (L5), and (CR) gives
$\nf(A^i\,W^i)=S\,A^{i+1}\,B^{i+1}$ and $\nf(b_i\,w_i)=\nf(B^i\,W^i)$ for $i<k$. Also
$a_k\,w_k\conv A^k\,W^k$ and $b_k\,w_k\conv B^k\,W^k$. (a) $W^i\psq B^i\,W^i$ and (RS). (b) At a
literal entry $a_{i-1}\,w_{i-1}=S\,a_i\,b_i$, so $b_i=w_{i-1}$; otherwise $b_i=y\,a_i$ is an
application, so all its reducts are (L3(e)). (c) $A^0=a_0$, $W^0=W$.
\end{proof}

So the move of a normal pair, with $k$, the new state and $N$, is computed from the pair by the
recursion above. The \emph{run} of $(A,W)$ is the sequence of states $(A_n,W_n)$ with tooth normal
forms $N_n$; by L5, $\rest^n(A\,W)\conv A_n\,W_n$ and the teeth of $(A_n,W_n)$ are convertible with
those of $A\,W$ from $n+1$ on. A move is a \emph{pop} if $k=0$ and \emph{raising} if $k\ge1$. A pop
keeps the heap and replaces the stack by a proper subterm; so raising moves recur. A raising move
replaces $W$ by $W^k$ with $W\psq W^k$. After two raising moves the heap has two proper right
suffixes, so it is neither $S$ nor $T$; after one it is not $S$.

\begin{lemma}\label{lem:last}
In a move of kind \textup{(B)}, if entry $i\ge1$ is not literal, then $w_{i-1}\psq b_i$.
\end{lemma}

\begin{proof}
The head run from $a_{i-1}\,w_{i-1}$ to $S\,a_i\,b_i$ has at least one step, and its frames before
$S\,a_i\,b_i$ have at least three items. The last item of $a_{i-1}\,w_{i-1}$ is $w_{i-1}$. A head step
from $S\,u_1\cdots u_m$ keeps the last item if $m\ge4$ and replaces it by $u_2\,u_3$ if $m=3$. So the
last item of each of these frames is $E[w_{i-1}]$ for a right context $E$, and in the frame
$S\,S\,y\,a_i$ before $S\,a_i\,b_i$ it is $a_i$. Hence $b_i=y\,E[w_{i-1}]$.
\end{proof}

\begin{lemma}\label{lem:reach}
Let $p\in\Ter$ have no hnf and only normalizing teeth. For some $n$, $\rest^n(p)\conv A\,W$ for a
normal pair $(A,W)$.
\end{lemma}

\begin{proof}
For a term $f=f^-\,u$ put $\|f\|=\sz{f^-}$. At a move of kind (A), $\rest(f)=f^-$ and
$\|\rest(f)\|<\|f\|$. At a move of kind (B) with every entry literal, $f^-=S\,a_0\,b_0$ and
$a_{i-1}=S\,a_i$, so $a_k=\rest(f)^-$ is a proper subterm of $f^-$. So some move $n$ of $p$ has a
non-literal entry; let $i$ be the last. Then $b_i=y\,a_i$ (also for $i=0$, by L3(a)) and
$a_i=S^{k-i}\,a_k$, so, as in Lemma~\ref{lem:pair}, $a_k$ and $w_k$ are subterms of the tooth
$b_k\,w_k$ and normalize. Put $A=\nf(a_k)$, $W=\nf(w_k)$: $A\,W\conv a_k\,w_k=\rest^n(p)$ has no
hnf (L5).
\end{proof}

\section{Normal forms of applications}

\begin{lemma}\label{lem:h2}
Let $B$, $V$ be normal with $B\,V$ normalizing. If $B$ has at most one item, $B\,V$ is normal. If
$\nf(B\,V)$ has one item, then $B=S$.
\end{lemma}

\begin{proof}
The first claim is clear. If $B=S\,b$, then $B\,V=S\,b\,V$ is normal with two items. If
$B=S\,x\,y$, then $B\,V$ has three items, and every later frame of its head run has at least two
(L3(a)); so its hnf has at least two items, and by L5 so has every term convertible with it, among
them $\nf(B\,V)$.
\end{proof}

\begin{lemma}\label{lem:fa}
Let $x=S\,x_1\,x_2$ be normal with two items, $y$ normal and $x\,y$ normalizing. Then
$\nf(x\,y)=S\,n_1\,n_2$ with $y\sqs n_1$, and $n_1=y$ iff $x_1=S$.
\end{lemma}

\begin{proof}
$x\,y\to(x_1\,y)\,(x_2\,y)$, so $g=\nf(x_1\,y)$ and $h=\nf(x_2\,y)$ exist (Lemma~\ref{lem:sub}),
$N=\nf(x\,y)=\nf(g\,h)$, $y\psq h$ and $h\psq N$ (RS); $g\ne S$ (L3(e)). Induction on the depth of
$y$ in $N$. If $g$ has one item, then $x_1=S$, $g=S\,y$ (Lemma~\ref{lem:h2}) and $N=S\,y\,h$. If $g$
has two items, then $x_1\ne S$, $h$ lies at a smaller depth in $N$ than $y$, and the induction
hypothesis for $(g,h)$ gives $N=S\,n_1\,n_2$ with $h\sqs n_1$; so $y\psq n_1$.
\end{proof}

\begin{lemma}\label{lem:tee}
Let $x$ be normal with two items, $y$ normal and $x\,y$ normalizing.
\textup{(a)} If $\nf(x\,y)=S\,c\,(S\,c)$, then $y=S$, or $x=T$ and $c=y$.
\textup{(b)} If $\nf(x\,y)=S\,c\,(S\,c\,(S\,c))$, then $y\in\{S,T\}$, or $x=S\,S\,T$ and $c=y$.
\end{lemma}

\begin{proof}
Let $x=S\,x_1\,x_2$, and $g$, $h$, $N$ as in the proof of Lemma~\ref{lem:fa}; induction on the
depth of $y$ in $N$. (a) Let $y\ne S$. If $g$ has one item, $x_1=S$ and $N=S\,y\,h$, so $c=y$ and
$h=\nf(x_2\,y)=S\,y$, whence $x_2=S$ (Lemma~\ref{lem:h2}) and $x=T$. If $g$ has two items, $h\ne S$
lies at a smaller depth, so by induction $g=T$. But $x_1=S$ makes $g$ one-item, $x_1=S\,e$ gives
$S\,e\,y=T$ and $y=S$, and a two-item $x_1$ gives $y\sqs S$ (Lemma~\ref{lem:fa}); each is excluded.
(b) Let $y\notin\{S,T\}$. If $g$ has one item, then $x_1=S$, $c=y$ and $h=\nf(x_2\,y)=S\,y\,(S\,y)$.
Here $x_2$ has two items, since $S\,y$ and $S\,e\,y$ differ from $S\,y\,(S\,y)$; so $x_2=T$ by (a),
and $x=S\,S\,T$. If $g$ has two items, then $h\notin\{S,T\}$, as $y\psq h$ and $S$ is the only
proper right suffix of $T$; by induction $g=S\,S\,T$. Now $x_1=S$, $x_1=S\,e$ and a two-item $x_1$
give a one-item $g$, $y=T$, and $y\sqs S$ respectively.
\end{proof}

A closed term $W$ is \emph{self-dead} if $W\,(t\,W)$ has no hnf for every closed $t$.

\begin{lemma}\label{lem:K}
\textup{(a)} For $c\in\Kc$ and normal $V$, $c\,V$ normalizes: $\nf(c\,V)=c\,V$ if $c$ has at most
one item, and $\nf(S\,S\,b\,V)=S\,V\,\nf(b\,V)$.
\textup{(b)} Let $d$ be closed and $\mathcal X$ a set of closed terms such that $d\,x$ has no hnf
for $x\in\mathcal X$ and $t\,d\in\mathcal X$ for every closed $t$. If $\alpha$ is normal,
$x_1,\dots,x_r\in\mathcal X$ and $\alpha\,d\,x_1\cdots x_r$ normalizes, then $r=0$ and $\alpha\in\Kc$,
or $r=1$ and $\alpha=S$.
\textup{(c)} If $W$ is self-dead, $c$ is normal and $c\,W$ normalizes, then $c\in\Kc$.
\end{lemma}

\begin{proof}
(a) $S\,S\,b\,V\to S\,V\,(b\,V)$; induction. (b) Induction on $\sz\alpha$. If $\alpha=S$ and $r\ge2$,
$S\,d\,x_1\,x_2\cdots\to d\,x_2\,(x_1\,x_2)\cdots$, which has no hnf (L3(d)). If $\alpha=S\,\alpha'$
and $r\ge1$, the head step gives $\alpha'\,x_1\,(d\,x_1)\cdots$, with a subterm without an hnf; so
$r=0$ and $\alpha\in\Kc$. If
$\alpha=S\,\alpha'\,\alpha''$, the head step gives $\alpha'\,d\,(\alpha''\,d)\,x_1\cdots x_r$, which
normalizes and has $r+1$ arguments from $\mathcal X$ after $d$; by induction $r=0$ and
$\alpha'=S$, so $S\,d\,(\alpha''\,d)$ normalizes, $\alpha''\,d$ does, $\alpha''\in\Kc$ by induction,
and $\alpha=S\,S\,\alpha''\in\Kc$. (c) Take $d=W$, $\mathcal X=\{t\,W\}$, $r=0$.
\end{proof}

\section{Moves over a self-dead heap}

\begin{lemma}\label{lem:dh}
Let $W$ be normal and self-dead.
\textup{(a)} If $(A,W)$ is a normal pair whose first tooth normalizes, then $B^0\in\Kc$, and either
the move is a pop, or $A^0=T$, $k=1$, the new state is $(W,\nf(B^0\,W))$ and
$N=S\,W\,\nf(B^0\,W)$.
\textup{(b)} \textup{(Landing.)} If $(C,H)$ is a normal pair with $H\ne S$ and first tooth
normalizing, whose move is raising with new heap $S\,W$, then its new state is $(W,S\,W)$.
\end{lemma}

\begin{proof}
(a) $B^0\,W$ normalizes (its normal form is $N$ if $k=0$ and $W^1$ otherwise), so $B^0\in\Kc$
(Lemma~\ref{lem:K}(c)). Let $k\ge1$. Then $W^1=\nf(B^0\,W)$ exists and $W\,W^1$ has no hnf, as
$W\,(B^0\,W)$ has none and reduces to it (L3(d)). Also $A^0\,W$ normalizes
(Lemma~\ref{lem:pair}), so $A^0\in\Kc$, and $A^0\ne S$, since $S\,S\,B^0\,W\to S\,W\,(B^0\,W)$.
If $A^0=S\,A^1$ has one item, then $B^1=W$, and $B^1\,W^1=W\,W^1$, whose normal form is $N$ or
$W^2$, has no hnf. So $A^0=S\,S\,b$ with $b\in\Kc$, and $\nf(A^0\,W)=S\,W\,\nf(b\,W)$: $A^1=W$,
$B^1=\nf(b\,W)$. As $A^1\,W^1=W\,W^1$ has no hnf, $k=1$. If $b=S\,b'$ or $b=S\,S\,b''$, then
$B^1\,W^1$ is $S\,b'\,W\,W^1$ or $S\,W\,\nf(b''\,W)\,W^1$, and its head step gives a term
containing $W\,W^1$; so $N$ would not exist. Hence $b=S$, $A^0=T$, $B^1=S\,W$ and $N=S\,W\,W^1$.

(b) In the move of $(C,H)$, $W^k=\nf(B^{k-1}\,W^{k-1})=S\,W$, so $B^{k-1}=S$ and $W^{k-1}=W$
(Lemma~\ref{lem:h2}). If $k\ge2$, entry $k-1$ is literal (Lemma~\ref{lem:pair}(b)), so
$W^{k-2}=S$; but $W^{k-2}$ is $H\ne S$ or has $H$ as a proper right suffix. So $k=1$ and $H=W$, and
(a) applies to $(C,W)$.
\end{proof}

\begin{lemma}\label{lem:alpha}
Let $\omega$ be normal and self-dead, and let $(\omega,V)$ and $(C,\omega)$ be normal pairs with
$\omega\psq V$. Their first teeth do not normalize to the same normal form.
\end{lemma}

\begin{proof}
Suppose both normalize to $N$. By Lemma~\ref{lem:dh}(a), $(C,\omega)$ pops some $B\in\Kc$ with
$N=\nf(B\,\omega)$, or makes the move (a) with $N=S\,\omega\,\nf(B\,\omega)$, $B\in\Kc$. In the
move of $(\omega,V)$ let $k$, $B^k$ and $W'=W^k$ be as in Lemma~\ref{lem:pair}; then $V\sqs W'\psq N$,
so $\omega\psq W'$ and $\sz{W'}>\sz\omega$. If $C$ pops a $B$ with at most one item, the proper right
suffixes of $N=B\,\omega$ are $\omega$ and its right suffixes, all smaller than $W'$. So
$N=S\,\omega\,V'$ with $V'=\nf(b\,\omega)$ for some $b\in\Kc$: $b=B$ in the second case, and if $C$
pops a two-item $B\in\Kc$, then $B=S\,S\,b$ and Lemma~\ref{lem:K}(a) applies. If $B^k$ had two items, Lemma~\ref{lem:fa} would give $W'\sqs\omega$; so
$B^k\,W'=N$, $B^k=S\,\omega$ and $W'=V'$. Let $\omega=S\,A^0\,B^0$.

If $k=0$, $B^k=B^0$ is smaller than $S\,\omega$. If $k\ge1$ and entry $k$ is not literal, then
$w_{k-1}\psq b_k$ (Lemma~\ref{lem:last}) and $W^{k-1}\psq B^k=S\,\omega$ by (RS), so
$V=W^0\sqs W^{k-1}\sqs\omega$, against $\sz V>\sz\omega$. If $k\ge1$ and entry $k$ is literal,
$W^{k-1}=B^k=S\,\omega$, and $V\sqs S\,\omega$ with $\sz V>\sz\omega$ gives $V=S\,\omega=W^{k-1}$;
as $W^0\psq W^1\psq\cdots$, $k=1$ and $V'=W^1=\nf(B^0\,(S\,\omega))$, with $S\,\omega\psq V'$. By
Lemma~\ref{lem:K}(a), $V'=S\,\omega\,(\cdots(S\,\omega\,(c\,\omega)))$ with $u\ge0$ factors
$S\,\omega$ and $c$ of at most one item; its right suffixes larger than $\omega$ are the $u+1$
terms obtained by deleting factors, and the one-item term $S\,\omega$ is one of them only if $c=S$
and $u\ge1$. So $V'=S\,\omega\,R$ for some $R$. If $B^0$ has at most one item, $B^0\,(S\,\omega)=V'$
and $B^0=S\,\omega$, larger than the subterm $B^0$ of $\omega$ can be. If $B^0$ has two items,
Lemma~\ref{lem:fa} gives $S\,\omega\sqs\omega$. All cases are contradictory.
\end{proof}

\section{The counter grammar}\label{sec:grammar}

Let $B=\{0,\dots,37\}$ with the table $\beta\colon B\times B\to B$ of~\citep{s1}; the \emph{block}
of a term is $\fold(S)=0$, $\fold(t\,u)=\beta(\fold t,\fold u)$. We use
\begin{itemize}
\item[(M2)] \citep[Proposition~4 and Theorem~9]{s1}, with L3(c): $t$ has no hnf iff $\fold(t)=36$.
\end{itemize}
Let $\Self=\{w\in B:\beta(w,\beta(\kappa,w))=36\text{ for all }\kappa\in B\}$. If
$\fold(W)\in\Self$, then $W$ is self-dead, since $\fold(W\,(t\,W))=36$. Put $\Theta_1=\{18,20,27\}$,
$\Theta_2=\{16,31,33,34\}$, $\Theta=\Theta_1\cup\Theta_2$, $CB=\{0,1,2,4,8,17\}$ and
$IB=\{12,22,30\}$.

\begin{lemma}\label{lem:tables}
$\fold(X)=16$ and $\fold(U)=33$. Moreover:
\begin{itemize}
\item[T16] \textup{(a)} $\beta(0,\theta)\in\Theta_1$ for $\theta\in\Theta$;
  $\beta(\beta(0,\rho),\kappa)\in\Theta_2$ for $\rho\in\Theta_2$, $\kappa\in B$;
  $\beta(\beta(0,\rho),\kappa)\in\Theta_2$ for $\rho\in\Theta_1$, $\kappa\in\Theta\cup\{0\}$.
  \textup{(b)} $\beta(\rho,\kappa)=36$ for $\rho\in\Theta_2$, $\kappa\in\Theta$.
  \textup{(c)} $\Theta_2\subseteq\Self$.
\item[T17] \textup{(a)} $\beta(0,\gamma)\in CB$ for $\gamma\in CB$;
  $\beta(\beta(0,\gamma),\kappa)\in IB$ for $\gamma\in CB$, $\kappa\in\Theta$.
  \textup{(b)} $\beta(\beta(0,\rho),\kappa)\in\Theta_2$ for $\rho\in\Theta_1$, $\kappa\in CB\cup IB$.
  \textup{(c)} $\beta(\rho,\kappa)=36$ for $\rho\in\Theta_2$, $\kappa\in IB$.
\end{itemize}
\end{lemma}

\begin{proof}
Finite computations on the table, checked by kernel evaluation in Lean (Data availability).
\end{proof}

All members of $\Hc$ and $\Ic$ are normal, since each rule applies $S$ to at most two normal terms.
Examples: $U_m\in\Hc_2$ by (H2); for $Y=X$ or $Y=S\,X\,c$ with $c$ normal, $Y\in\Hc_2$, $S\,Y$ and
$S\,(S\,Y)$ are in $\Hc$, and so are the terms $G_e$, $J_e$ of Section~\ref{sec:families}.

\begin{lemma}\label{lem:ch}
\textup{(a)} $S$, $T$, the combs and the members of $\Ic$ are not in $\Hc$. The one-item members
of $\Hc$ are the terms $S\,H$ with $H\in\Hc$.
\textup{(b)} Let $Z=S\,Z^0\,B^0\in\Hc_2\setminus\{X,U\}$. Then $Z^0\in\Hc$. If $Z^0$ has one item,
then $B^0\in\Hc\cup\Ic\cup\Cc$, and $B^0\in\Cc\setminus\{S\}$ only if $Z^0=S\,(S\,H)$ with
$H\in\Hc$. In any case $Z^0=S^m\,P_0$ for some $m\ge0$ and $P_0\in\Hc_2$; put $\pop(Z)=P_0$.
\end{lemma}

\begin{proof}
(H1) makes one-item terms $S\,H$, the other rules make $X$, $U$ (first item $T$) and two-item
terms $S\,P\,Q$ with $P\in\Hc$ (from (H2)--(H4); $P$ has one item exactly in (H3), (H4)), and a term
determines its items. No rule makes $S$. So $T$, whose first item is $S$, is not in $\Hc$;
$c_{j+1}=S\,c_j\in\Hc$ only if $c_j\in\Hc$, so no comb is; and the first item of $S\,c\,Q\in\Ic$ is a
comb, neither $T$ nor a member. This gives (a) and the first claims of (b); for the clause on
$\Cc\setminus\{S\}$, note that these combs have one item, so they are neither in $\Hc$ nor in $\Ic$,
whose members have two, and only (H4) admits them. Stripping the leading $S$'s from $Z^0$ by (a)
stays in $\Hc$, never reaches $S$, and ends at a two-item member.
\end{proof}

\begin{lemma}\label{lem:dead}
\textup{(a)} $\fold(t)\in\Theta_1$ for one-item $t\in\Hc$, $\in\Theta_2$ for $t\in\Hc_2$, $\in IB$
for $t\in\Ic$, and $\in CB$ for $t\in\Cc$.
\textup{(b)} $Z\,t$ has no hnf for $Z\in\Hc_2$ and $t\in\Hc\cup\Ic$; so every member of $\Gc$ is a
normal pair.
\textup{(c)} Every $Z\in\Hc_2$ is self-dead.
\end{lemma}

\begin{proof}
(a) Induction on size, by Lemma~\ref{lem:tables}: (H0) by the folds of $X$, $U$; (H1) by T16(a);
(H2) by T16(a); (H3) by T16(a) for $Q\in\Hc\cup\{S\}$ and T17(b) for $Q\in\Ic$; (H4) by T17(b), as
$S\,P'$ is a one-item member; (I) and the combs by T17(a). (b) T16(b), T17(c) and (M2). (c) T16(c).
\end{proof}

\begin{lemma}\label{lem:inert}
Let $I=S\,c_j\,Q\in\Ic$ and $y\in\Hc$ with $I\,y$ normalizing. Then $Q=S\,Q'$, and with
$u=S\,Q'\,y$, either $j=0$ and $\nf(I\,y)=S\,y\,u$, or $j=1$, $y=S\,y'$ and
$\nf(I\,y)=S\,u\,(S\,y'\,u)$. In both cases $\nf(I\,y)\in\Hc_2$.
\end{lemma}

\begin{proof}
$I\,y\to c_j\,y\,(Q\,y)$. If $Q$ has two items, $Q\,y$ has no hnf (Lemma~\ref{lem:dead}(b)), and
$I\,y$ does not normalize. So
$Q=S\,Q'$ with $Q'\in\Hc$, and $u=Q\,y$ is normal and in $\Hc_2$ by (H2) or (H3). If $j=0$,
$S\,y\,u$ is normal and in $\Hc_2$ by (H2) or (H3). If $j\ge1$, $c_j\,y\,u\to c_{j-1}\,u\,(y\,u)$, and
$y\,u$ has no hnf unless $y=S\,y'$; then $y\,u=S\,y'\,u\in\Hc_2$. For $j=1$ the term
$S\,u\,(S\,y'\,u)$ is normal and in $\Hc_2$ by (H2). For $j\ge2$,
$c_{j-1}\,u\,(y\,u)\to c_{j-2}\,(y\,u)\,(u\,(y\,u))$ contains $u\,(y\,u)$, which has no hnf.
\end{proof}

\begin{theorem}\label{thm:closure}
Let $(Z,W)\in\Gc$ have a normalizing first tooth. Then its move is one of the following, and its
new state is in $\Gc$:
\begin{itemize}
\item[(i)] $Z\notin\{X,U\}$: with $Z^0=S^m\,\pop(Z)$, $k=m$ and the new state is
  $(\pop(Z),W^m)$; and $m=0$ if $W$ has two items;
\item[(ii)] $Z=X$, $W$ with two items: $k=1$, new state $(W,S\,W)$, $N=S\,W\,(S\,W)$;
\item[(iii)] $Z=U$, $W$ with two items: $k=1$, new state $(W,D(W))$, $N=S\,W\,D(W)$;
\item[(iv)] $Z=X$, $W=S\,V$ with $V$ two-item: $k=2$, new state $(V,D(S\,V))$.
\end{itemize}
\end{theorem}

\begin{proof}
We use Lemma~\ref{lem:pair} throughout. First, for $x\in\Hc\cup\Ic\cup\Cc$ and $y\in\Hc\cup\Ic$ with
$x\,y$ normalizing:
\begin{itemize}
\item[($*$)] $x\notin\Hc_2$ (Lemma~\ref{lem:dead}(b)); if $x=S\,x'\in\Hc$, then $\nf(x\,y)=S\,x'\,y\in\Hc_2$
  by (H2) or (H3); if $x\in\Ic$ and $y\in\Hc$, then $\nf(x\,y)\in\Hc_2$ (Lemma~\ref{lem:inert}); if
  $x=c_j$ and $y\in\Hc$, then $x\,y$ is $S\,y\in\Hc$ or $S\,c_{j-1}\,y\in\Ic$.
\end{itemize}
(i) Let $Z^0=S^m\,P_0$ with $P_0=\pop(Z)$ (Lemma~\ref{lem:ch}). For $i<m$, $A^i=S^{m-i}\,P_0$ and
$A^i\,W^i$ is normal with two items, so $A^{i+1}=S^{m-i-1}\,P_0$ and $B^{i+1}=W^i$. Hence $k\ge m$,
$W^1=\nf(B^0\,W)$ and $W^{i+1}=\nf(W^{i-1}\,W^i)$ for $1\le i<m$, and these exist. We show
$W^i\in\Hc\cup\Ic$ for $i\le m$, with $W^i\in\Ic$ only if $i=1$ and $m\ge2$. If $m\ge1$, then
$B^0\in\Hc\cup\Ic\cup\Cc$, and by ($*$) $W^1\in\Hc\cup\Ic$, inert only if $B^0\in\Cc\setminus\{S\}$,
which forces $m\ge2$ (Lemma~\ref{lem:ch}(b)). For $1\le i<m$, $W^{i-1}$ is not $S$ (it is $W\in\Hc$
or the normal form of an application) and not in $\Hc_2$, since $\nf(W^{i-1}\,W^i)$ exists; if it
is inert, $i-1=1$ and $W^i=W^2\in\Hc$. By ($*$), $W^{i+1}\in\Hc_2$. So $W^m\in\Hc$, $P_0\,W^m$ has
no hnf (Lemma~\ref{lem:dead}(b)), $k=m$, and the new state is in $\Gc$. If $W$ has two items and
$m\ge1$, the normal form of $W\,W^1$ is the tooth or $W^2$, but $W\,W^1$ has no hnf; so $m=0$.

(ii), (iii) $Z=S\,T\,b$ with $b=S$ resp.\ $T$. As $T\,W\to S\,W\,(S\,W)$, normal, we get $A^1=W$,
$B^1=S\,W$, $W^1=\nf(b\,W)$, which is $S\,W\in\Hc$ resp.\ $D(W)\in\Hc_2$ (H1, H2). $W$ is self-dead
(Lemma~\ref{lem:dead}(c)), so $A^1\,W^1$ has no hnf, $k=1$, and $N=\nf(S\,W\,W^1)=S\,W\,W^1$.

(iv) $\nf(T\,(S\,V))=S\,(S\,V)\,(S\,(S\,V))$, so $A^1=S\,V$, $B^1=W^1=S\,(S\,V)$. $A^1\,W^1$ is normal,
so $A^2=V$, $B^2=W^1$, $W^2=\nf(B^1\,W^1)=D(S\,V)$, which is in $\Hc$ by (H3). $V\,D(S\,V)$ has no hnf,
so $k=2$.

Finally $Z=X$ over $W=S\,W'$ with $W'=S\,V'$, and $Z=U$ over a one-item $W=S\,V$, do not occur.
In the first, as in (ii), $A^1=W$, $B^1=W^1=S\,W$; then $A^2=W'$, $B^2=S\,W$, $W^2=D(W)$; then
$A^3=V'$, $B^3=D(W)$, $W^3=S\,W\,D(W)$, and $\nf(B^3\,W^3)$, the tooth or $W^4$, would be the normal
form of $D(W)\,(S\,W\,D(W))$, a member of $\Hc_2$ applied to a member of $\Hc$ (H3). In the second,
$A^1=W$, $B^1=S\,W$, $W^1=D(W)$; $A^2=V$, $B^2=D(W)$, $W^2=S\,W\,D(W)$, and $\nf(B^2\,W^2)$ fails
in the same way.
\end{proof}

\begin{corollary}\label{cor:recur}
Let $(Z,W)\in\Gc$ have only normalizing teeth, and let $(A_n,W_n)$ be its run. Every
$(A_n,W_n)$ is in $\Gc$, and for every $n_0$ there is $n>n_0$ with $A_{n-1}\in\{X,U\}$ and $W_{n-1}\in\Hc_2$, so that move $n$
is \textup{(ii)} or \textup{(iii)}.
\end{corollary}

\begin{proof}
By Theorem~\ref{thm:closure}. Moves (i) make the stack smaller, so moves (ii)--(iv) recur. After
(iv) the heap has two items; over it moves (i) keep the heap and (iv) cannot occur, so the next move
other than (i) is (ii) or (iii).
\end{proof}

\section{Proof of Theorem~\ref{thm:main}}

\begin{proposition}\label{prop:meet}
Let $P_0\in\Gc$ and $Q_0$ be normal pairs whose teeth all normalize, with the same normal forms
$N_n$. Their runs reach equal states at some index.
\end{proposition}

\begin{proof}
Let $(A_n,W_n)$ and $(A^\circ_n,W^\circ_n)$ be the runs of $P_0$ and $Q_0$. Choose $n_0$ such that
two of the moves $1,\dots,n_0$ of $Q_0$ are raising. Then $W^\circ_i\notin\{S,T\}$ for $i\ge n_0$, and from the second raising move of $Q_0$ on, the
heap before each move is not $S$.
By Corollary~\ref{cor:recur} take $n>n_0$ such that move $n$ of $P_0$ is (ii) or (iii): with
$W=W_{n-1}\in\Hc_2$, $(A_n,W_n)=(W,V)$ with $V=S\,W$ resp.\ $V=D(W)$, and $N_n=S\,W\,V$. $W$ is
self-dead (Lemma~\ref{lem:dead}(c)), and so is $D(W)\in\Hc_2$ (H2). In move $n$ of $Q_0$ let $k$,
$B=B^k$, $(A_q,W_q)=(A^\circ_n,W^\circ_n)$ be as in Lemma~\ref{lem:pair}, so $N_n=\nf(B\,W_q)$. If
$B$ has at most one item, $B\,W_q=S\,W\,V$, so $B=S\,W$ and $W_q=V$. If $B$ has two items,
Lemma~\ref{lem:fa} gives $W_q\sqs W$. Four cases remain.

\emph{$W_q=W$.} $(W,V)$ and $(A_q,W)$ are normal pairs with $W\psq V$, and their first teeth both
normalize to $N_{n+1}$, against Lemma~\ref{lem:alpha}.

\emph{$W_q\psq W$.} Then $B$ has two items, and by Lemma~\ref{lem:tee}, applied to
$N_n=S\,W\,(S\,W)$ resp.\ $S\,W\,(S\,W\,(S\,W))$, $W_q\in\{S,T\}$ or $W_q=W$: both are excluded.

\emph{$W_q=V=S\,W$.} If $k\ge1$, Lemma~\ref{lem:dh}(b), with $H=W^\circ_{n-1}\ne S$, gives
$(A_q,W_q)=(W,S\,W)=(A_n,W_n)$. If $k=0$, $Q_0$ pops $B=S\,W$ over $W^\circ_{n-1}=S\,W$, so
$A^\circ_{n-1}=S\,A_q\,(S\,W)$. Let $j<n$ be the last raising move of $Q_0$. Its new heap is
$W^\circ_j=W^\circ_{n-1}=S\,W$, and the heap before it is not $S$, so $A^\circ_j=W$ by
Lemma~\ref{lem:dh}(b). The pops $j+1,\dots,n-1$ take proper subterms, so $A^\circ_{n-1}$ is a
subterm of $W$; but it is larger than $W$.

\emph{$W_q=V=D(W)$.} Both runs now stand on the self-dead heap $V$. Let $m>n$ be least such that
move $m$ of $P_0$ or of $Q_0$ is raising; it exists, since raising moves recur. Before $m$ both runs
pop and keep the heap $V$. By Lemma~\ref{lem:dh}(a) a raising move at $m$ leads to $(V,V')$ with
$V'=\nf(b\,V)$ and tooth $N_m=S\,V\,V'$. If both runs raise, their new heaps are both the second
item of $N_m$, and the states are equal. If one raises and the other pops, they stand at $(V,V')$
and $(C,V)$ with $V\psq V'$, and their teeth $m+1$ both normalize to $N_{m+1}$, against
Lemma~\ref{lem:alpha}.
\end{proof}

\begin{proof}[Proof of Theorem~\ref{thm:main}]
By L5, $\rest^e(q)$ has no hnf and its teeth are convertible with $\tau_{e+i}(p)$ and with the teeth
of $Z\,W$, so they normalize. By Lemma~\ref{lem:reach}, $\rest^{e+d}(q)\conv Q_0$ for some $d$ and
some normal pair $Q_0$. Let $P_0$ be the state of the run of $(Z,W)$ at index $d$; it is in $\Gc$
(Corollary~\ref{cor:recur}) and $\rest^{e+d}(p)\conv P_0$. The teeth of $P_0$ and $Q_0$ are
convertible with $\tau_{e+d+i}(p)\conv\tau_{e+d+i}(q)$, so they have the same normal forms (CR). By
Proposition~\ref{prop:meet} the runs of $P_0$ and $Q_0$ reach equal states at some index $n$, so
$\rest^{e+d+n}(p)\conv\rest^{e+d+n}(q)$ (L5), and L6 gives $p\conv q$.
\end{proof}

\section{The two families}\label{sec:families}

\begin{lemma}\label{lem:four}
Let $a,b,c,w,W$ be closed.
\textup{(a)} If $a\,w$ has no hnf, the move of $S\,a\,c\,w$ has tooth $c\,w$ and rest $a\,w$.
\textup{(b)} If $W\,(b\,W)$ has no hnf, the move of $S\,T\,b\,W$ has tooth $S\,W\,(b\,W)$ and rest
$W\,(b\,W)$.
\textup{(c)} If $a\,(c\,w)$ has no hnf, the move of $S\,(S\,a)\,c\,w$ has tooth $w\,(c\,w)$ and rest
$a\,(c\,w)$.
\textup{(d)} If $a\,D(S\,a)$ has no hnf, the move of $X\,(S\,a)$ has tooth $S\,(S\,a)\,D(S\,a)$ and rest
$a\,D(S\,a)$.
In each case the term has no hnf.
\end{lemma}

\begin{proof}
The head runs
\begin{gather*}
S\,a\,c\,w\to a\,w\,(c\,w),\qquad
S\,(S\,a)\,c\,w\to S\,a\,w\,(c\,w)\to a\,(c\,w)\,(w\,(c\,w)),\\
S\,T\,b\,W\to T\,W\,(b\,W)\to S\,W\,(S\,W)\,(b\,W)\to W\,(b\,W)\,(S\,W\,(b\,W)),\\
X\,v\to T\,v\,(S\,v)\to S\,v\,(S\,v)\,(S\,v)\to S\,a\,(S\,v)\,D(v)\to a\,D(v)\,(S\,v\,D(v)),
\end{gather*}
where $v=S\,a$, reach the claimed rest applied to further arguments, so the terms have no hnf
(L3(c),(d)). L4(B) gives the
moves: (a) $a_0=a$, $b_0=c$, $w_0=w$, and $a_0\,w_0$ has no hnf, so $k=0$. (b) $a_0=T$, $b_0=b$,
$w_0=W$; $T\,W\to S\,W\,(S\,W)$ gives $a_1=W$, $b_1=S\,W$, $w_1=b\,W$; $k=1$. (c) $a_0\,w_0=S\,a\,w$
is literal: $a_1=a$, $b_1=w$, $w_1=c\,w$; $k=1$. (d) $a_1=v$, $b_1=w_1=S\,v$; $a_1\,w_1=S\,a\,(S\,v)$
is literal, so $a_2=a$, $b_2=S\,v$, $w_2=D(v)$; $k=2$.
\end{proof}

\begin{proof}[Proof of Corollary~\ref{cor:U}]
$(U_j,U_m)\in\Gc$, and $U_j\,U_m$ has no hnf (Lemma~\ref{lem:dead}(b)). For $j\ge1$,
$U_j=S\,U_{j-1}\,(S\,U_{j-1})$, and Lemma~\ref{lem:four}(a) gives the normal tooth
$S\,U_{j-1}\,U_m$ and the rest $U_{j-1}\,U_m$. For $j=0$, Lemma~\ref{lem:four}(b) with $b=T$
($U_m$ is self-dead) gives the tooth $S\,U_m\,(T\,U_m)\red S\,U_m\,U_{m+1}$ and the rest
$U_m\,(T\,U_m)\conv U_m\,U_{m+1}$. By induction and L5 every tooth normalizes. Apply
Theorem~\ref{thm:main}.
\end{proof}

For the second family let $G_0=D(S\,Y)$, $J_{e+1}=S\,(S\,Y)\,(S\,G_e)$ and
$G_{e+1}=S\,G_e\,(S\,J_{e+1})$. By (H1)--(H3), $Y$, $G_e$, $J_e\in\Hc_2$ and $S\,Y$, $S\,G_e$,
$S\,J_e\in\Hc$. The run of $X\,Y$ is literally the following, each line by
Lemma~\ref{lem:four}; the lines marked $\ell$ occur only for $Y\ne X$, and for $Y=X$ the term
$Y\,V$ in the other lines reads $X\,V$.
\begin{center}\small
$\begin{array}{@{}llll@{}}
\text{move} & \text{term}\to\text{rest} & \text{tooth} & \text{by}\\[2pt]
(1) & X\,Y\to Y\,(S\,Y) & S\,Y\,(S\,Y) & \text{(b)}, b=S\\
(1\ell) & Y\,(S\,Y)\to X\,(S\,Y) & c\,(S\,Y) & \text{(a)}\\
(2) & X\,(S\,Y)\to Y\,G_0 & S\,(S\,Y)\,G_0 & \text{(d)}\\
(\ell_e) & Y\,G_e\to X\,G_e & c\,G_e & \text{(a)}\\
(\alpha_e) & X\,G_e\to G_e\,(S\,G_e) & S\,G_e\,(S\,G_e) & \text{(b)}, b=S\\
(\beta_{e,i}),\ i=e,\dots,1 & G_i\,(S\,G_e)\to G_{i-1}\,(S\,G_e) & S\,J_i\,(S\,G_e) & \text{(a)}\\
(\gamma_e) & G_0\,(S\,G_e)\to Y\,J_{e+1} & S\,G_e\,J_{e+1} & \text{(c)}\\
(\ell'_e) & Y\,J_{e+1}\to X\,J_{e+1} & c\,J_{e+1} & \text{(a)}\\
(\delta_e) & X\,J_{e+1}\to J_{e+1}\,(S\,J_{e+1}) & S\,J_{e+1}\,(S\,J_{e+1}) & \text{(b)}, b=S\\
(\varepsilon_e) & J_{e+1}\,(S\,J_{e+1})\to Y\,G_{e+1} & S\,J_{e+1}\,G_{e+1} & \text{(c)}
\end{array}$
\end{center}
After (2) comes phase $0$; phase $e\ge0$ is $(\ell_e)$, $(\alpha_e)$, $(\beta_{e,e}),\dots,(\beta_{e,1})$,
$(\gamma_e)$, $(\ell'_e)$, $(\delta_e)$, $(\varepsilon_e)$, followed by phase $e+1$. In $(\gamma_e)$,
$G_0=S\,(S\,Y)\,(S\,(S\,Y))$ and $S\,(S\,Y)\,(S\,G_e)=J_{e+1}$; in $(\varepsilon_e)$,
$S\,G_e\,(S\,J_{e+1})=G_{e+1}$. Every hypothesis of Lemma~\ref{lem:four} used is that a member of
$\Hc_2$ applied to a member of $\Hc$ has no hnf (Lemma~\ref{lem:dead}(b)); in (b), $W\,(S\,W)$ with
$W\in\{Y,G_e,J_{e+1}\}$.

\begin{proof}[Proof of Corollary~\ref{cor:X}]
$(X,Y)\in\Gc$. Every tooth in the table is normal except the teeth $c\,V$ of the lines $\ell$. If
$c\in\Kc$ they normalize (Lemma~\ref{lem:K}(a)), and Theorem~\ref{thm:main} applies. If $c\notin\Kc$,
the tooth $c\,G_0$ of $(\ell_0)$ does not normalize, since $G_0\in\Hc_2$ is self-dead
(Lemmas~\ref{lem:dead}(c) and~\ref{lem:K}(c)).
\end{proof}

\section{An open question}

Call a move of a normal pair a \emph{refill} if some entry is not literal. An all-literal move
replaces the stack $A=S\,A^0\,B^0$ by $A^k$ with $A^0=S^k\,A^k$, a proper subterm; so refills recur.

\begin{conjecture}\label{conj:entry}
The run of every normal pair whose teeth all normalize reaches $\Gc$; more sharply, its state after
its second refill is in $\Gc$.
\end{conjecture}

\noindent Since refills recur, the sharper form implies the first. With Lemma~\ref{lem:reach} and
Theorem~\ref{thm:main}, the first gives CT against every $p$ whose teeth normalize. A proof must
control what the first refills do to arbitrary initial stacks and heaps, before the run is a
counter.

\section*{Data availability}
\sloppy
The proofs in this paper are not mechanized. Lemma~\ref{lem:tables} and (M2) are. In
\texttt{lean-ski}~\citep{leanski} (Lean~4~\citep{lean4}), version 1.20.0 (merge commit
\texttt{20d8271}), module \texttt{Ski/Lock1Tables.lean}, with only \texttt{propext} and
\texttt{Quot.sound}: T16 and T17 are \texttt{Lock1Tables.T16} and \texttt{Lock1Tables.T17}, proved by
\texttt{decide +kernel} ($\Theta_1$, $\Theta_2$, $CB$, $IB$ are \texttt{TH1}, \texttt{TH2},
\texttt{CB}, \texttt{IB}), with the term forms \texttt{T16\_dead} and \texttt{T17\_dead}; the folds of
$X$ and $U$ are \texttt{fold\_X} and \texttt{fold\_U}; (M2) is \texttt{hnfless\_iff}, from
\texttt{HnfSound.sound} and \texttt{HnfS.no\_hnf\_of\_bot}; $\Self$ and the self-deadness of terms
with block in $\Self$ are \texttt{mem\_SELF} and \texttt{self\_dead}. (CR), (Std) and (F2) are as
in~\citep{s11}: \texttt{church\_rosser}, \texttt{wsteps\_split} and \texttt{Full.run\_total}.

\bibliographystyle{plainnat}
\bibliography{refs}
\end{document}
